\subsection{域 R 上的向量空间与代数的拓扑}

设 \(\mathbf{E}\) 是 \(n\) 维的一个向量空间（域 \(\mathbf{R}\) 上的）；若 \((\boldsymbol{a}_i)_{1 \leqslant i \leqslant n}\) 是 \(\mathbf{E}\) 的一个基，则每一点 \(\boldsymbol{x} \in \mathbf{E}\) 都可唯一地写成形式 \(\boldsymbol{x} = \sum_{i=1}^{n} x_i \boldsymbol{a}_i\) ，其中诸 \(x_i\) 是实数；因此映射 \((x_i) \to \sum_{i=1}^{n} x_i \boldsymbol{a}_i\) 是 \(\mathbf{R}^n\) 到 \(\mathbf{E}\) 上的一个双射线性映射。如果 \(\mathbf{E}\) 经此映射被赋以 \(\mathbf{R}^n\) 的拓扑，那么 \(\mathbf{E}\) 就有一个与其加法群结构相容的拓扑，而且映射 \((t, \boldsymbol{x}) \to tx\) （\(\mathbf{R} \times \mathbf{E}\) 到 \(\mathbf{E}\) 中的）关于这个拓扑是连续的。这个拓扑不依赖于 \(\mathbf{E}\) 中所选取的基；因为若 \((\boldsymbol{a}')\) 是 \(\mathbf{E}\) 的另一个基，且 \(x = \sum_{i=1}^{n} x_i' \boldsymbol{a}_i' = \sum_{i=1}^{n} x_i \boldsymbol{a}_i\) ，则映射 \((x_i) \to (x_i')\) （\(\mathbf{R}^n\) 到其自身的）是线性的，因此是一个同胚。

这个事实使人猜测，这样定义在 E 上的拓扑应当能够不借助 E 的基而得到刻画。事实上，我们以后将看到，它是 E 上使函数 x - y（在 E × E 上）与 tx（在 R × E 上）为连续的唯一的豪斯多夫拓扑。

现在若 A 是域 R 上有限秩 n 的一个代数，则上面在 A 上的拓扑（把 A 看作 R 上 n 维向量空间）不仅与 A 的加法群结构相容，而且与其环结构相容。这是下述更一般结果的一个推论：

命题 5. 设 E, F, G 是域 R 上三个有限维向量空间。则 E × F 到 G 中的每个双线性映射（*）f 是连续的。

我们可以设 \(E = R^{m}\) ，\(F = R^{n}\) ，\(G = R^{p}\) ；只需证明 \(R^{p}\) 中 \(f(x, y)\) 的坐标是 \((x, y) \in E \times F\) 的连续函数（第 I 章，§ 4，第 1 小节，命题 1）。换句话说，只需证明每个双线性形式 g 在 \(E \times F\) 上连续；而这是直接的，因为 \(g(x, y)\) 是 x 与 y 的坐标的多项式。

